% doc/jva_plots.tex - generated by drawJVA.C from rootfiles/JVA_v2.root (tag _v2); do not edit
\providecommand{\plotdir}{../plots/}
\providecommand{\pt}{p_{\mathrm{T}}}

\subsection{Missing transverse momentum}
The missing transverse momentum of every usable vertex, $\vec{p}_{\mathrm{T}}^{\,\mathrm{miss}}=-\sum w_i\,\vec{p}_{\mathrm{T},i}$ over the candidates of its per-vertex event under each method, for all vertices, the primary vertex and the others (figure~\ref{fig:jva_met_dist}); its width against $N_{\mathrm{PU}}$ and against $S_v$, the scalar $\pt$ sum of the vertex-resolved part (figure~\ref{fig:jva_met_width}); and the residual against $-\sum\vec{p}_{\mathrm{T}}$ of the generated particles of the owner interaction with $|\eta|<5$, the imbalance a perfect detector would see. The rms is the in-axis rms; the core is a Gaussian iterated to convergence in $\mu\pm2\sigma$, failed below 200 entries, without convergence, when it is wider than the distribution or off its median by more than $\sigma/2$. At the primary vertex \texttt{dup} and \texttt{lv} are the same event, the standard reconstruction. The sample has 25000 crossings at $\langle N_{\mathrm{PU}}\rangle=45$ with 34.3 usable vertices per crossing, 91\% of them owned by an interaction (within 0.05\,cm). Anti-$k_{\mathrm{T}}$ $R=0.4$ from 1\,GeV; candidates beyond $|\eta|=2.5$ are vertex-blind; track candidates need 0.5\,GeV in the cone; the optimiser uses $\sigma_v^2=1^2+1^2\,S_v$\,GeV$^2$ and the prices $\tau_{\mathrm{trk}}=4$, $\tau_{\mathrm{all}}=12$; the oracles accept a vertex within 0.2\,cm. Jets are corrected with \texttt{text/jec\_v1.txt}. PUPPI weights from \texttt{puppi.h} version \texttt{genseed-1.0}.
\begin{figure}[htbp]\centering
\includegraphics[width=0.32\textwidth]{\plotdir jva_metxy_all.pdf}\hfil
\includegraphics[width=0.32\textwidth]{\plotdir jva_metxy_pv.pdf}\hfil
\includegraphics[width=0.32\textwidth]{\plotdir jva_metxy_pu.pdf}\\
\includegraphics[width=0.32\textwidth]{\plotdir jva_metabs_all.pdf}\hfil
\includegraphics[width=0.32\textwidth]{\plotdir jva_metabs_pv.pdf}\hfil
\includegraphics[width=0.32\textwidth]{\plotdir jva_metabs_pu.pdf}\\
\includegraphics[width=0.32\textwidth]{\plotdir jva_metres_all.pdf}\hfil
\includegraphics[width=0.32\textwidth]{\plotdir jva_metres_pv.pdf}\hfil
\includegraphics[width=0.32\textwidth]{\plotdir jva_metres_pu.pdf}
\caption{Missing transverse momentum per vertex, unit-normalised: first row the $x$ and $y$ components, second row the magnitude, third row the residual of the components against the owner interaction's true imbalance (owned vertices), for all vertices, the primary vertex and the other vertices (left to right), for \texttt{dup} (grey), \texttt{lv} (blue), \texttt{trk} (green), \texttt{jva} (red), \texttt{ojet} (open orange), \texttt{opart} (open magenta), \texttt{none} (open black and dashed). The legend quotes the rms and the core $\sigma$ (or the mean magnitude). Over all vertices the rms of a component is \texttt{dup} 14.9\,GeV, \texttt{lv} 7.2\,GeV, \texttt{jva} 6.9\,GeV, \texttt{opart} 6.9\,GeV, \texttt{none} 7.0\,GeV. The core width of the residual against the truth is \texttt{dup} 13.0\,GeV, \texttt{lv} 4.9\,GeV, \texttt{jva} 4.9\,GeV, \texttt{opart} 4.9\,GeV, \texttt{none} 4.9\,GeV.}
\label{fig:jva_met_dist}
\end{figure}

\begin{figure}[htbp]\centering
\includegraphics[width=0.32\textwidth]{\plotdir jva_metw_npu.pdf}\hfil
\includegraphics[width=0.32\textwidth]{\plotdir jva_metw_s.pdf}\hfil
\includegraphics[width=0.32\textwidth]{\plotdir jva_metresw_s.pdf}
\caption{Width of the missing transverse momentum of all vertices: the components against $N_{\mathrm{PU}}$ and against $S_v$, and the residual against the owner interaction's true imbalance against $S_v$ (owned vertices), with the noise model of the optimiser, $\sqrt{\sigma_0^2+\sigma_K^2 S_v}$ (dashed). Upper panels the rms, lower panels the core $\sigma$; methods as in figure~\ref{fig:jva_met_dist}. The model describes the vertex-resolved part against the owner's true central recoil; the residual here is against the true imbalance over $|\eta|<5$ and so also holds the part of the owner's forward recoil that no forward cluster carries (PUPPI keeps little of it), which is why it lies above the model at large $S_v$. At $S_v$ in 30--40\,GeV the model gives 6.0\,GeV, the residual core of \texttt{opart} is 7.4\,GeV and of \texttt{jva} 7.4\,GeV.}
\label{fig:jva_met_width}
\end{figure}


\subsection{The forward spike}
Jets of the per-vertex events of the owned vertices under each method, against the pure generated jets of the owner interaction (filled once per owned vertex): $\mathrm{d}N/\mathrm{d}|\eta|$ per owned vertex above three thresholds in $\pt^{\mathrm{raw}}$ (figure~\ref{fig:jva_spike}), the $\pt^{\mathrm{corr}}$ spectra in three regions (figure~\ref{fig:jva_spectra}), and the purity of the jets above 5\,GeV: the fraction of a jet's linked generated $\pt$ that comes from the owner interaction, and the share of its dominant interaction (figure~\ref{fig:jva_purity}). A jet made of several interactions has a dominant share well below 1; the spike at $2.5<|\eta|<3.0$ is made of them.
\begin{figure}[htbp]\centering
\includegraphics[width=0.32\textwidth]{\plotdir jva_spike_pt05.pdf}\hfil
\includegraphics[width=0.32\textwidth]{\plotdir jva_spike_pt10.pdf}\hfil
\includegraphics[width=0.32\textwidth]{\plotdir jva_spike_pt20.pdf}
\caption{Jets per owned vertex against $|\eta|$ above 5, 10 and 20\,GeV (left to right; reconstructed jets in $\pt^{\mathrm{raw}}$): the pure generated jets of the owner interaction (black line) and the reconstructed jets of the owned vertex's event under \texttt{dup} (grey), \texttt{lv} (blue), \texttt{trk} (green), \texttt{jva} (red), \texttt{ojet} (open orange), \texttt{opart} (open magenta), \texttt{none} (open black and dashed), with the ratio to the generated jets below (on a logarithmic scale where the methods span more than a factor 20); the dotted lines mark $2.5<|\eta|<3.0$. At $2.5<|\eta|<3.0$ above 10\,GeV the reconstructed jets per owned vertex are 49.61 (\texttt{dup}), 1.53 (\texttt{lv}), 1.80 (\texttt{trk}), 0.10 (\texttt{jva}), 0.24 (\texttt{ojet}), 0.15 (\texttt{opart}), 0.00 (\texttt{none}) times the generated ones.}
\label{fig:jva_spike}
\end{figure}

\begin{figure}[htbp]\centering
\includegraphics[width=0.32\textwidth]{\plotdir jva_spectra_r0.pdf}\hfil
\includegraphics[width=0.32\textwidth]{\plotdir jva_spectra_r1.pdf}\hfil
\includegraphics[width=0.32\textwidth]{\plotdir jva_spectra_r2.pdf}
\caption{$\pt$ spectra per owned vertex of the pure generated jets of the owner interaction and of the reconstructed jets ($\pt^{\mathrm{corr}}$, corrected to the genseed link scale) under each method, at $|\eta|<2.5$, $2.5<|\eta|<3.0$ and $3.0<|\eta|<5.0$ (left to right), with the ratio to the generated spectrum below (on a logarithmic scale where the methods span more than a factor 20). At $2.5<|\eta|<3.0$ and 10--12\,GeV reco/gen is 45.92 (\texttt{dup}), 1.43 (\texttt{lv}), 1.75 (\texttt{trk}), 0.06 (\texttt{jva}), 0.20 (\texttt{ojet}), 0.18 (\texttt{opart}), not defined (\texttt{none}).}
\label{fig:jva_spectra}
\end{figure}

\begin{figure}[htbp]\centering
\includegraphics[width=0.48\textwidth]{\plotdir jva_own.pdf}\hfil
\includegraphics[width=0.48\textwidth]{\plotdir jva_dom_r0.pdf}\\
\includegraphics[width=0.48\textwidth]{\plotdir jva_dom_r1.pdf}\hfil
\includegraphics[width=0.48\textwidth]{\plotdir jva_dom_r2.pdf}
\caption{Purity of the jets above 5\,GeV at the owned vertices: the fraction of the linked generated $\pt$ that comes from the owner interaction against $|\eta|$ (the profile), and the distribution of the dominant interaction's share of the linked $\pt$ at $|\eta|<2.5$, $2.5<|\eta|<3.0$ and $3.0<|\eta|<5.0$ (a share of exactly 1 is in the last bin). The mean own fraction at $2.5<|\eta|<3.0$ is \texttt{dup} 0.03, \texttt{lv} 0.13, \texttt{trk} 0.17, \texttt{jva} 0.13, \texttt{ojet} 0.61, \texttt{opart} 0.79, \texttt{none} 0.35. Fraction of the jets at $2.5<|\eta|<3.0$ with a dominant share below 0.4 / above 0.6: \texttt{dup} 0.73 / 0.06, \texttt{lv} 0.70 / 0.07, \texttt{trk} 0.70 / 0.07, \texttt{jva} 0.70 / 0.10, \texttt{ojet} 0.01 / 0.53, \texttt{opart} 0.01 / 0.85, \texttt{none} 0.25 / 0.25.}
\label{fig:jva_purity}
\end{figure}


\subsection{Forward-cluster assignment}
Every cluster of the vertex-blind candidates ($|\eta|>2.5$, anti-$k_{\mathrm{T}}$ $R=0.4$) is classified by the dominant interaction of its linked generated $\pt$: \emph{single} when that share exceeds 0.5 and the interaction has a usable vertex within 0.2\,cm, \emph{combination} otherwise (unlinked clusters included). Its fate under each method that assigns whole clusters is then \emph{right} (a single cluster sent to the vertex of its interaction), \emph{wrong} (sent elsewhere), \emph{single nulled}, \emph{combination nulled} (the right thing to do with a combination) or \emph{combination kept}; \emph{correct} is right plus combination nulled. The clusters with track candidates (charged candidates of a vertex above 0.5\,GeV in the cone at $|\eta|<2.5$) and those without are shown apart: without tracks the optimiser can place a cluster at any vertex, at the price $\tau_{\mathrm{all}}=12$.
\begin{figure}[htbp]\centering
\includegraphics[width=0.32\textwidth]{\plotdir jva_fwd_clpt.pdf}\hfil
\includegraphics[width=0.32\textwidth]{\plotdir jva_fwd_share.pdf}\hfil
\includegraphics[width=0.32\textwidth]{\plotdir jva_fwd_ncand.pdf}
\caption{The forward clusters: their $\pt$ spectrum per crossing with and without track candidates and the fraction with tracks, the fraction whose dominant interaction has more than 0.5, more than 0.6 and less than 0.4 of the linked $\pt$ (an unlinked cluster has share 0), and the number of track candidates. Above 5\,GeV there are 5.03 clusters per crossing, 52\% of them with track candidates. At 10--12\,GeV 0.21 of the clusters are single (share $>0.5$) and 0.65 have a share below 0.4. A cluster has 1.72 track candidates on average; 33\% have none.}
\label{fig:jva_fwd}
\end{figure}

\begin{figure}[htbp]\centering
\includegraphics[width=0.24\textwidth]{\plotdir jva_cat_lv_trk.pdf}\hfil
\includegraphics[width=0.24\textwidth]{\plotdir jva_cat_trk_trk.pdf}\hfil
\includegraphics[width=0.24\textwidth]{\plotdir jva_cat_jva_trk.pdf}\hfil
\includegraphics[width=0.24\textwidth]{\plotdir jva_cat_ojet_trk.pdf}\\
\includegraphics[width=0.24\textwidth]{\plotdir jva_cat_lv_notrk.pdf}\hfil
\includegraphics[width=0.24\textwidth]{\plotdir jva_cat_trk_notrk.pdf}\hfil
\includegraphics[width=0.24\textwidth]{\plotdir jva_cat_jva_notrk.pdf}\hfil
\includegraphics[width=0.24\textwidth]{\plotdir jva_cat_ojet_notrk.pdf}
\caption{Outcome fractions of the forward clusters against the cluster $\pt$ for \texttt{lv}, \texttt{trk}, \texttt{jva} and \texttt{ojet} (left to right), for the clusters with track candidates (first row) and without (second row): single cluster at the right vertex (green), at a wrong vertex (red), nulled (orange); combination nulled (blue) or kept (magenta). Above 5\,GeV the correct fraction is \texttt{lv} 0.06, \texttt{trk} 0.05, \texttt{jva} 0.74, \texttt{ojet} 1.00, \texttt{none} 0.76 (clusters with and without tracks together), where \texttt{none}, dropping every cluster, scores the fraction of combinations. The outcomes add up to the cluster spectrum in every bin (largest relative difference 0).}
\label{fig:jva_outcome}
\end{figure}

\begin{figure}[htbp]\centering
\includegraphics[width=0.48\textwidth]{\plotdir jva_correct_trk.pdf}\hfil
\includegraphics[width=0.48\textwidth]{\plotdir jva_correct_notrk.pdf}
\caption{The correct fraction (right vertex plus combination nulled) of the forward clusters against the cluster $\pt$, with track candidates (left) and without (right), per method, and the reference of dropping every cluster (\texttt{none}, black dashed), which scores the fraction of combinations: a method that places clusters helps only where it lies above that line. Above 5\,GeV with tracks: \texttt{lv} 0.037, \texttt{trk} 0.031, \texttt{jva} 0.818, \texttt{ojet} 1.000; \texttt{none} 0.857. Above 5\,GeV without tracks: \texttt{lv} 0.081, \texttt{trk} 0.081, \texttt{jva} 0.659, \texttt{ojet} 1.000; \texttt{none} 0.660.}
\label{fig:jva_correct}
\end{figure}


\subsection{Relative response from every vertex}
Dijets are selected from the per-vertex event of every usable vertex: the two leading jets in $\pt^{\mathrm{corr}}$ above 3\,GeV at $|\eta|<5.191$, $\Delta\phi>2.7$, a tag at $|\eta|<1.305$ and the other jet as the probe (both assignments at weight 1/2 when both are in the barrel), $\alpha=\pt^{(3)}/\pt^{\mathrm{avg}}$. $R_{\mathrm{DB}}=(1+\langle A_{\mathrm{DB}}\rangle)/(1-\langle A_{\mathrm{DB}}\rangle)$ with $A_{\mathrm{DB}}=(\pt^{\mathrm{probe}}-\pt^{\mathrm{tag}})/(\pt^{\mathrm{probe}}+\pt^{\mathrm{tag}})$, and the same for $A_{\mathrm{MPF}}=\vec{p}_{\mathrm{T}}^{\,\mathrm{miss,T1}}\cdot\hat{p}_{\mathrm{T}}^{\,\mathrm{tag}}/(2\pt^{\mathrm{avg}})$ (type-1 correction above 3\,GeV). The closure reference $r_{\mathrm{true}}$ is the median of the per-event ratio of the probe's to the tag's $\pt^{\mathrm{corr}}/\pt^{\mathrm{linkz}}$, where $\pt^{\mathrm{linkz}}$ is the linked generated $\pt$ of the jet. A mean or a median needs 50 effective dijets.
\begin{figure}[htbp]\centering
\includegraphics[width=0.32\textwidth]{\plotdir jva_l2_db_pt05.pdf}\hfil
\includegraphics[width=0.32\textwidth]{\plotdir jva_l2_db_pt08.pdf}\hfil
\includegraphics[width=0.32\textwidth]{\plotdir jva_l2_db_pt12.pdf}
\caption{$R_{\mathrm{DB}}$ against the probe $|\eta|$ for the dijets of all vertices with $\alpha<0.3$, at $\pt^{\mathrm{avg}}$ 5--6, 8--10 and 12--15\,GeV (left to right, top to bottom; of the five L2Res bins, those in which some method has two probe bins with 50 effective dijets), per method (markers, slightly shifted in $|\eta|$ for legibility), with the truth relative response $r_{\mathrm{true}}$ of the same dijets (lines, same colour); below, the closure $R/r_{\mathrm{true}}$, with $\pm2\%$ dotted. The mean $|R_{\mathrm{DB}}/r_{\mathrm{true}}-1|$ over the probe bins where both are measured (their number in parentheses) is, at 5--6\,GeV \texttt{lv} 0.089 (11), \texttt{trk} 0.144 (13), \texttt{jva} 0.090 (11), \texttt{ojet} 0.095 (11), \texttt{opart} 0.091 (11), \texttt{none} 0.090 (11); at 8--10\,GeV \texttt{lv} 0.127 (11), \texttt{trk} 0.264 (12), \texttt{jva} 0.124 (11), \texttt{ojet} 0.125 (11), \texttt{opart} 0.134 (11), \texttt{none} 0.124 (11); at 12--15\,GeV \texttt{lv} 0.223 (2), \texttt{trk} 0.317 (1), \texttt{jva} 0.231 (2), \texttt{ojet} 0.243 (2), \texttt{opart} 0.295 (1), \texttt{none} 0.231 (2).}
\label{fig:jva_l2_db}
\end{figure}

\begin{figure}[htbp]\centering
\includegraphics[width=0.32\textwidth]{\plotdir jva_l2_mpf_pt05.pdf}\hfil
\includegraphics[width=0.32\textwidth]{\plotdir jva_l2_mpf_pt08.pdf}\hfil
\includegraphics[width=0.32\textwidth]{\plotdir jva_l2_mpf_pt12.pdf}
\caption{$R_{\mathrm{MPF}}$ against the probe $|\eta|$ for the dijets of all vertices with $\alpha<0.3$, at $\pt^{\mathrm{avg}}$ 5--6, 8--10 and 12--15\,GeV (left to right, top to bottom; of the five L2Res bins, those in which some method has two probe bins with 50 effective dijets), per method (markers, slightly shifted in $|\eta|$ for legibility), with the truth relative response $r_{\mathrm{true}}$ of the same dijets (lines, same colour); below, the closure $R/r_{\mathrm{true}}$, with $\pm2\%$ dotted. The mean $|R_{\mathrm{MPF}}/r_{\mathrm{true}}-1|$ over the probe bins where both are measured (their number in parentheses) is, at 5--6\,GeV \texttt{lv} 0.059 (11), \texttt{trk} 0.104 (13), \texttt{jva} 0.059 (11), \texttt{ojet} 0.062 (11), \texttt{opart} 0.060 (11), \texttt{none} 0.059 (11); at 8--10\,GeV \texttt{lv} 0.106 (11), \texttt{trk} 0.239 (12), \texttt{jva} 0.103 (11), \texttt{ojet} 0.104 (11), \texttt{opart} 0.115 (11), \texttt{none} 0.103 (11); at 12--15\,GeV \texttt{lv} 0.250 (2), \texttt{trk} 0.290 (1), \texttt{jva} 0.257 (2), \texttt{ojet} 0.263 (2), \texttt{opart} 0.284 (1), \texttt{none} 0.257 (2).}
\label{fig:jva_l2_mpf}
\end{figure}

\begin{figure}[htbp]\centering
\includegraphics[width=0.32\textwidth]{\plotdir jva_l2_db_pt05_a1.pdf}\hfil
\includegraphics[width=0.32\textwidth]{\plotdir jva_l2_db_pt08_a1.pdf}\hfil
\includegraphics[width=0.32\textwidth]{\plotdir jva_l2_db_pt12_a1.pdf}\\
\includegraphics[width=0.32\textwidth]{\plotdir jva_l2_db_pt18_a1.pdf}
\caption{$R_{\mathrm{DB}}$ against the probe $|\eta|$ for the dijets of all vertices without the $\alpha$ cut (the $\alpha<0.3$ selection is in figure~\ref{fig:jva_l2_db} and figure~\ref{fig:jva_l2_mpf}), at $\pt^{\mathrm{avg}}$ 5--6, 8--10, 12--15 and 18--21\,GeV (left to right, top to bottom; of the five L2Res bins, those in which some method has two probe bins with 50 effective dijets), per method (markers, slightly shifted in $|\eta|$ for legibility), with the truth relative response $r_{\mathrm{true}}$ of the same dijets (lines, same colour); below, the closure $R/r_{\mathrm{true}}$, with $\pm2\%$ dotted. The mean $|R_{\mathrm{DB}}/r_{\mathrm{true}}-1|$ over the probe bins where both are measured (their number in parentheses) is, at 5--6\,GeV \texttt{dup} 0.257 (2), \texttt{lv} 0.061 (11), \texttt{trk} 0.102 (13), \texttt{jva} 0.061 (11), \texttt{ojet} 0.061 (11), \texttt{opart} 0.058 (11), \texttt{none} 0.061 (11); at 8--10\,GeV \texttt{dup} 0.103 (15), \texttt{lv} 0.030 (11), \texttt{trk} 0.084 (13), \texttt{jva} 0.070 (13), \texttt{ojet} 0.040 (13), \texttt{opart} 0.050 (13), \texttt{none} 0.031 (11); at 12--15\,GeV \texttt{dup} 0.097 (17), \texttt{lv} 0.045 (12), \texttt{trk} 0.057 (13), \texttt{jva} 0.074 (13), \texttt{ojet} 0.056 (13), \texttt{opart} 0.050 (13), \texttt{none} 0.044 (11); at 18--21\,GeV \texttt{dup} 0.079 (14), \texttt{lv} 0.068 (12), \texttt{trk} 0.071 (13), \texttt{jva} 0.075 (11), \texttt{ojet} 0.087 (12), \texttt{opart} 0.069 (11), \texttt{none} 0.071 (10).}
\label{fig:jva_l2_db_a1}
\end{figure}

\begin{figure}[htbp]\centering
\includegraphics[width=0.32\textwidth]{\plotdir jva_l2_mpf_pt05_a1.pdf}\hfil
\includegraphics[width=0.32\textwidth]{\plotdir jva_l2_mpf_pt08_a1.pdf}\hfil
\includegraphics[width=0.32\textwidth]{\plotdir jva_l2_mpf_pt12_a1.pdf}\\
\includegraphics[width=0.32\textwidth]{\plotdir jva_l2_mpf_pt18_a1.pdf}
\caption{$R_{\mathrm{MPF}}$ against the probe $|\eta|$ for the dijets of all vertices without the $\alpha$ cut (the $\alpha<0.3$ selection is in figure~\ref{fig:jva_l2_db} and figure~\ref{fig:jva_l2_mpf}), at $\pt^{\mathrm{avg}}$ 5--6, 8--10, 12--15 and 18--21\,GeV (left to right, top to bottom; of the five L2Res bins, those in which some method has two probe bins with 50 effective dijets), per method (markers, slightly shifted in $|\eta|$ for legibility), with the truth relative response $r_{\mathrm{true}}$ of the same dijets (lines, same colour); below, the closure $R/r_{\mathrm{true}}$, with $\pm2\%$ dotted. The mean $|R_{\mathrm{MPF}}/r_{\mathrm{true}}-1|$ over the probe bins where both are measured (their number in parentheses) is, at 5--6\,GeV \texttt{dup} 0.034 (2), \texttt{lv} 0.038 (11), \texttt{trk} 0.084 (13), \texttt{jva} 0.039 (11), \texttt{ojet} 0.040 (11), \texttt{opart} 0.041 (11), \texttt{none} 0.038 (11); at 8--10\,GeV \texttt{dup} 0.074 (15), \texttt{lv} 0.033 (11), \texttt{trk} 0.087 (13), \texttt{jva} 0.069 (12), \texttt{ojet} 0.030 (13), \texttt{opart} 0.076 (13), \texttt{none} 0.032 (11); at 12--15\,GeV \texttt{dup} 0.095 (17), \texttt{lv} 0.063 (12), \texttt{trk} 0.060 (13), \texttt{jva} 0.120 (13), \texttt{ojet} 0.055 (13), \texttt{opart} 0.069 (12), \texttt{none} 0.056 (11); at 18--21\,GeV \texttt{dup} 0.080 (14), \texttt{lv} 0.098 (12), \texttt{trk} 0.087 (13), \texttt{jva} 0.113 (11), \texttt{ojet} 0.081 (12), \texttt{opart} 0.083 (11), \texttt{none} 0.084 (10).}
\label{fig:jva_l2_mpf_a1}
\end{figure}


\subsection{Jet energy resolution from dijets}
Dijets whose two jets fall into the same $|\eta|$ bin, with $\alpha<0.3$, are projected on the bisector axis $\hat{n}=(\hat{u}_1-\hat{u}_2)/|\hat{u}_1-\hat{u}_2|$, along which both jets make the same angle, and on its normal: $(\vec{p}_{\mathrm{T},1}+\vec{p}_{\mathrm{T},2})\cdot\hat{n}/\pt^{\mathrm{avg}}$ for the DB bisector and $\vec{p}_{\mathrm{T}}^{\,\mathrm{miss,T1}}\cdot\hat{n}/\pt^{\mathrm{avg}}$ for MPFX. The parallel width holds both jets' resolution and the imbalance along the axis, the perpendicular one the imbalance across it, so the resolution per jet is $\sigma=\sqrt{(\mathrm{RMS}_\parallel^2-\mathrm{RMS}_\perp^2)/2}$; a bin with $\mathrm{RMS}_\perp\geq\mathrm{RMS}_\parallel$ has no JER and counts as a failed bin. Each distribution is first mirror-symmetrised ($h(x)+h(-x)$: the sign of $\hat{n}$ follows the order of the two jets, and a $\pt$-ordered pair folds the parallel projection at zero); the width is then the core $\sigma$ of a Gaussian centred at zero, iterated in $\pm2\sigma$ (primary), or the plain in-axis rms. Each needs 200 dijets per $\pt^{\mathrm{avg}}$ bin, on a $\pt^{\mathrm{avg}}$ axis of 2 merged fine bins from 3\,GeV (3--5, 5--8, 8--12, 12--18\,GeV, \ldots). The truth is the width of $\pt^{\mathrm{corr}}/\pt^{\mathrm{linkz}}$ of both jets of the same dijets, like with like: $\sigma_{\mathrm{core}}/\mu_{\mathrm{core}}$ under the response protocol for the core-fit JER, rms/mean for the plain-rms JER. The \emph{reach} is the lowest $\pt^{\mathrm{avg}}$ bin with $|\mathrm{JER}/\mathrm{JER}_{\mathrm{true}}-1|<0.1$ whose next two measured bins are not significantly outside that band ($|\mathrm{JER}/\mathrm{JER}_{\mathrm{true}}-1|<0.1+2\sigma$, and no failed bin). Besides the 18 $|\eta|$ bins, the same analysis is made in four regions of whole bins (barrel, endcap, transition, HF), whose same-bin samples are the sums of their bins' (figure~\ref{fig:jva_jer_db_regions} and figure~\ref{fig:jva_jer_mpfx_regions}).
\begin{figure}[htbp]\centering
\includegraphics[width=0.48\textwidth]{\plotdir jva_jer_db_r0.pdf}\hfil
\includegraphics[width=0.48\textwidth]{\plotdir jva_jer_db_r1.pdf}\\
\includegraphics[width=0.48\textwidth]{\plotdir jva_jer_db_r2.pdf}
\caption{JER per jet from the DB bisector method against $\pt^{\mathrm{avg}}$ for the same-bin dijets of all vertices ($\alpha<0.3$) in four regions made of whole $|\eta|$ bins, $0<|\eta|<1.305$, $1.305<|\eta|<2.5$, $2.5<|\eta|<3.139$ and $3.139<|\eta|<5.191$ (left to right, top to bottom; the regions without 200 dijets in any $\pt^{\mathrm{avg}}$ bin are left out): core-fit JER per method (markers) and the truth $\sigma_{\mathrm{core}}/\mu_{\mathrm{core}}$ of $\pt^{\mathrm{corr}}/\pt^{\mathrm{linkz}}$ of the same jets (lines, same colour); below, the ratio with $\pm10\%$ dotted, and a cross on the bottom edge where $\mathrm{RMS}_\perp\geq\mathrm{RMS}_\parallel$.}
\label{fig:jva_jer_db_regions}
\end{figure}

\begin{figure}[htbp]\centering
\includegraphics[width=0.32\textwidth]{\plotdir jva_jer_db_e07.pdf}\hfil
\includegraphics[width=0.32\textwidth]{\plotdir jva_jer_db_e08.pdf}
\caption{JER per jet from the DB bisector method against $\pt^{\mathrm{avg}}$ for the same-bin dijets of all vertices ($\alpha<0.3$), per $|\eta|$ bin from $0<|\eta|<0.261$ to $1.93<|\eta|<2.172$ (left to right, top to bottom; the bins without 200 dijets in any $\pt^{\mathrm{avg}}$ bin are left out); markers, lines and the lower panels as in figure~\ref{fig:jva_jer_db_regions}.}
\label{fig:jva_jer_db_0}
\end{figure}

\begin{figure}[htbp]\centering
\includegraphics[width=0.32\textwidth]{\plotdir jva_jer_db_e12.pdf}
\caption{JER per jet from the DB bisector method against $\pt^{\mathrm{avg}}$ for the same-bin dijets of all vertices ($\alpha<0.3$), per $|\eta|$ bin from $2.172<|\eta|<2.322$ to $3.839<|\eta|<5.191$ (left to right, top to bottom; the bins without 200 dijets in any $\pt^{\mathrm{avg}}$ bin are left out); markers, lines and the lower panels as in figure~\ref{fig:jva_jer_db_regions}.}
\label{fig:jva_jer_db_1}
\end{figure}

\begin{figure}[htbp]\centering
\includegraphics[width=0.48\textwidth]{\plotdir jva_jer_mpfx_r0.pdf}\hfil
\includegraphics[width=0.48\textwidth]{\plotdir jva_jer_mpfx_r1.pdf}\\
\includegraphics[width=0.48\textwidth]{\plotdir jva_jer_mpfx_r2.pdf}
\caption{JER per jet from the MPFX method against $\pt^{\mathrm{avg}}$ for the same-bin dijets of all vertices ($\alpha<0.3$) in four regions made of whole $|\eta|$ bins, $0<|\eta|<1.305$, $1.305<|\eta|<2.5$, $2.5<|\eta|<3.139$ and $3.139<|\eta|<5.191$ (left to right, top to bottom; the regions without 200 dijets in any $\pt^{\mathrm{avg}}$ bin are left out): core-fit JER per method (markers) and the truth $\sigma_{\mathrm{core}}/\mu_{\mathrm{core}}$ of $\pt^{\mathrm{corr}}/\pt^{\mathrm{linkz}}$ of the same jets (lines, same colour); below, the ratio with $\pm10\%$ dotted, and a cross on the bottom edge where $\mathrm{RMS}_\perp\geq\mathrm{RMS}_\parallel$. Reach at $0<|\eta|<1.305$: \texttt{opart} none. Reach at $1.305<|\eta|<2.5$: \texttt{lv} none, \texttt{jva} none, \texttt{opart} none, \texttt{none} none.}
\label{fig:jva_jer_mpfx_regions}
\end{figure}

\begin{figure}[htbp]\centering
\includegraphics[width=0.32\textwidth]{\plotdir jva_jer_mpfx_e07.pdf}\hfil
\includegraphics[width=0.32\textwidth]{\plotdir jva_jer_mpfx_e08.pdf}
\caption{JER per jet from the MPFX method against $\pt^{\mathrm{avg}}$ for the same-bin dijets of all vertices ($\alpha<0.3$), per $|\eta|$ bin from $0<|\eta|<0.261$ to $1.93<|\eta|<2.172$ (left to right, top to bottom; the bins without 200 dijets in any $\pt^{\mathrm{avg}}$ bin are left out); markers, lines and the lower panels as in figure~\ref{fig:jva_jer_mpfx_regions}.}
\label{fig:jva_jer_mpfx_0}
\end{figure}

\begin{figure}[htbp]\centering
\includegraphics[width=0.32\textwidth]{\plotdir jva_jer_mpfx_e12.pdf}
\caption{JER per jet from the MPFX method against $\pt^{\mathrm{avg}}$ for the same-bin dijets of all vertices ($\alpha<0.3$), per $|\eta|$ bin from $2.172<|\eta|<2.322$ to $3.839<|\eta|<5.191$ (left to right, top to bottom; the bins without 200 dijets in any $\pt^{\mathrm{avg}}$ bin are left out); markers, lines and the lower panels as in figure~\ref{fig:jva_jer_mpfx_regions}.}
\label{fig:jva_jer_mpfx_1}
\end{figure}


\subsection{Summary}
The key number of each section per method is collected in table~\ref{tab:jva_summary} (vertex level) and table~\ref{tab:jva_summary_dijet} (dijets).
